\documentclass[12pt]{article}
\usepackage[utf8]{inputenc}
\usepackage{amsmath, amssymb, amsfonts, mathtools, array}
\usepackage{geometry}
\geometry{a4paper, margin=1in}
\title{Student Answer Script Transcription (High-Resolution Stitched Document)}
\author{Roll Number: 240101079}
\date{\today}
\begin{document}
\maketitle

\textbf{Q-1)} Space for answering Question-1 only.

\textbf{a)} For $S$ to be subspace of $M_{4\times4}(R)$, it must satisfy the following:

\textit{i)} $S$ must include zero matrix.

In zero matrix, as all the elements are zero,
the condition for $S$ is satisfied $(0+0+0=0)$ for fourth column.

\section*{Closure Under Addition}

Let $W_{1}$ and $W_{2}$ be two elements from $S$. Then $W_{1} + W_{2}$ must belong to $S$.

Each element in $W_{1} + W_{2}$ is the sum of respective elements of $W_{1}$ and $W_{2}$. So, the fourth column of $(W_{1} + W_{2})$ = Fourth col. of $W_{1}$ + Fourth col. of $W_{2}$, which can be further broken down to sum of first 3 columns of $W_{1}$ + sum of first 3 columns of $W_{2}$. By associativity of addition, $(W_{1}+W_{2})$ belongs to $S$.

\section*{Closure Under Scalar Multiplication}

If $W \in S$, then $tW$ must belong to $S$. As in row $\left[\begin{array}{ccc}a_{1} & a_{2} & a_{3}\\a_{4} & a_{5} & a_{6}\end{array}\right]$, $ta_{1}+ta_{2}+ta_{3}=t(a_{1}+a_{2}+a_{3})$ (distributive and commutative).

\end{document}
